% Part: first-order-logic
% Chapter: syntax
% Section: intro-syntax

\documentclass[../../../include/open-logic-section]{subfiles}

\begin{document}

% SEGMENT OLP-0150-S01
\olfileid{fol}{syn}{itx}

\olsection{அறிமுகம்}

முதல்தர தருக்கத்தின் கோட்பாட்டையும் மேல்கோட்பாட்டையும் வளர்க்க, அதன்
வெளிப்பாடுகளின் தொடரமைப்பையும் பொருண்மையியலையும் முதலில் வரையறுக்க வேண்டும்.
முதல்தர தருக்கத்தின் வெளிப்பாடுகள் சொற்களும் !!{formula}s உம் ஆகும். சொற்கள்
!!{variable}s, !!{constant}s மற்றும் !!{function}s ஆகியவற்றிலிருந்து
உருவாக்கப்படுகின்றன. !!^{formula}s, மறுபுறம், !!{predicate}s மற்றும் சொற்களுடன்
சேர்ந்து உருவாகின்றன (இவை மிகச் சிறிய, ``அணு'' !!{formula}s); பின்னர் அணு
!!{formula}s இலிருந்து தருக்கத் தொடுப்புக்குறிகளையும் அளவையடைகளையும் கொண்டு
மிகச் சிக்கலானவற்றை உருவாக்கலாம். உருவாக்க விதிகளை எழுதுவதற்குப் பல வழிகள்
உள்ளன; அவற்றில் ஒரு சாத்தியமான வழியை மட்டுமே தருகிறோம். பிற அமைப்புகள்
வேறு குறிகளைத் தேர்ந்தெடுக்கலாம், முதன்மையான தொடுப்புக்குறிகளின் வேறு
தொகுப்பைத் தேர்ந்தெடுக்கலாம், அடைப்புக்குறிகளை வேறுவிதமாகப் பயன்படுத்தலாம்
(அல்லது போலிஷ் குறியீட்டில் போல அவற்றைப் பயன்படுத்தாமலும் இருக்கலாம்).
எல்லா அணுகுமுறைகளுக்கும் பொதுவானது என்னவென்றால், உருவாக்க விதிகள் சொற்களின்
மற்றும் !!{formula}s இன் தொகுப்பை \emph{தொகுத்தறிதல் வழியாக} வரையறுக்கின்றன.
அது முறையாகச் செய்யப்பட்டால், ஒவ்வொரு வெளிப்பாடும் உருவாக்க விதிகளின்படி
அடிப்படையில் ஒரே ஒரு வழியில் மட்டுமே பெறப்படும். வெளிப்பாடுகள்
\emph{தனித்த வாசிப்புடையவை} ஆகும் என்ற இந்தத் தொகுத்தறிதல் வரையறை,
அதே முறையை---தொகுத்தறிதல் வரையறையைப் பயன்படுத்தி---அவற்றுக்கு
பொருள்களையும் வழங்க வழி செய்கிறது.

\end{document}
